\newcommand{\figuraPresionProfundidad}{%
\begin{figure}[htbp]
    \centering
    \begin{tikzpicture}[>=Latex,font=\small]
        \draw[very thick] (1.0,5.0)--(1.0,0.6)--(8.0,0.6)--(8.0,5.0);
        \fill[cyan!20] (1.03,0.63) rectangle (7.97,4.25);
        \draw[very thick,blue!65!black] (1.0,4.25)--(8.0,4.25);

        \fill[red!75] (4.5,1.35) circle (0.10);
        \node[red!75!black,anchor=west] at (4.68,1.35)
            {punto 1: $p_1$};
        \fill[blue!75] (4.5,3.55) circle (0.10);
        \node[blue!70!black,anchor=west] at (4.68,3.55)
            {punto 2: $p_2$};

        \draw[dashed,gray!70] (2.2,1.35)--(4.5,1.35);
        \draw[dashed,gray!70] (2.2,3.55)--(4.5,3.55);
        \draw[decorate,decoration={brace,amplitude=6pt},thick]
            (2.0,1.35)--node[left=6pt]{$h$}(2.0,3.55);

        \draw[ultra thick,green!45!black,-{Latex[length=4mm]}]
            (9.2,4.25)--(9.2,2.95);
        \node[green!40!black,anchor=west] at (9.38,3.55)
            {$\vec{g}=-g_0\vec{e}_z$};

        \node[font=\bfseries] at (4.5,5.35)
            {Presión y profundidad};
        \node[align=center] at (5.0,0.02)
            {$p_1=p_2+\rho g_0h$ \qquad $p_1>p_2$};
    \end{tikzpicture}
    \caption{La diferencia de presión entre dos puntos de un líquido depende
    únicamente de su separación vertical. El punto más profundo tiene mayor
    presión.}
    \label{fig:presion-profundidad}
\end{figure}%
}